% Project status, 2026-11-24. Plain beamer, no branding.
% Compile: pdflatex status_2026-11-24.tex (twice). Numbers from PLAN.md, PLAN_COSMO.md,
% MILESTONES.md as of 2026-11-24.
\documentclass[11pt,aspectratio=169,xcolor=table]{beamer}

\usepackage[utf8]{inputenc}
\usepackage[T1]{fontenc}
\usepackage{lmodern}
\usepackage{booktabs}
\usepackage{amssymb}
\usepackage{array}

\usetheme{default}
\setbeamertemplate{navigation symbols}{}
\setbeamertemplate{footline}{\leavevmode\hfill\insertframenumber\,/\,\inserttotalframenumber\hspace*{1em}\par\vskip0.5em}
\setbeamertemplate{itemize items}{\textbullet}
\setbeamerfont{frametitle}{series=\bfseries}
\setbeamercolor{frametitle}{fg=black}
\setbeamercolor{title}{fg=black}
\setbeamercolor{structure}{fg=black!75}

\definecolor{done}{gray}{0.55}
\definecolor{plan}{gray}{0.85}
\newcommand{\yes}{\textbf{yes}}
\newcommand{\no}{no}
\newcommand{\cdone}{\cellcolor{done}}
\newcommand{\cplan}{\cellcolor{plan}}
\newcommand{\take}[1]{\vfill\begin{block}{}#1\end{block}}
\setbeamercolor{block body}{bg=black!6}

\title{Unsupervised domain adaptation\\for precipitation downscaling}
\subtitle{State of the project}
\author{Filippo Quarenghi}
\date{24 November 2026}

\begin{document}

%----------------------------------------------------------------------
\begin{frame}[plain]
\titlepage
\end{frame}

%----------------------------------------------------------------------
\begin{frame}{The problem}
\begin{itemize}
  \item High-resolution rain observations are dense over Europe and North
        America, sparse over most of the tropics.
  \item Downscaling models are trained where observations exist and used
        where they do not. Moved to another region, they lose much of their
        skill.
  \item \textbf{Unsupervised domain adaptation (UDA)} proposes a cheap fix:
        during training, add the target region's \emph{inputs} (no rain
        observations) and make the model treat both regions alike.
  \item We test whether this works for precipitation downscaling, and when
        it can work at all.
\end{itemize}
\end{frame}

%----------------------------------------------------------------------
\begin{frame}{Research questions}
\begin{enumerate}
  \item Does the relation \emph{rain given the atmosphere}, $P(Y\mid X)$,
        differ between regions, or only the atmosphere, $P(X)$?
  \item How much of the transfer gap could be recovered at best, and can
        unlabelled target data find it?
  \item Do UDA methods do inside the network what they claim to do?
  \item Does the answer change with the model (deterministic vs.\
        generative)?
  \item Does a perfect-model setting, where the relation is shared by
        construction, change the answer?
\end{enumerate}
\take{Organising idea: UDA can only help if $P(Y\mid X)$ is the same in both
regions. That can be measured before any adaptation is trained.}
\end{frame}

%----------------------------------------------------------------------
\begin{frame}{Setting}
\small
\begin{tabular}{@{}lll@{}}
\toprule
 & \textbf{RainShift} & \textbf{COSMO perfect model} \\
\midrule
Input $\to$ target & ERA5 $\to$ IMERG & COSMO, coarse $\to$ 2.2\,km \\
Coarse rain input & ERA5 rain, another product & exact block mean of the target \\
Grid & $80^2 \to 200^2$ ($2.5\times$) & $20^2 \to 200^2$ ($10\times$) \\
Regions & Europe W, Horn of Africa, Melanesia & 5 European regions \\
 & (+ N.\ America E, W.\ Africa, no methods) & \\
Directions & 6 with methods (20 without) & 20 \\
\bottomrule
\end{tabular}

\medskip
\begin{itemize}
  \item \textbf{Models:} U-Net (deterministic); AFM, a flow-matching model
        that starts from the U-Net's mean; CFM, plain conditional flow
        matching (coded, not trained yet).
  \item \textbf{Methods:} FDA, AdaBN, DANN, MMD, DeepJDOT, a spectral loss.
        References: the \emph{oracle} (trained with target labels) and
        few-shot fine-tuning.
  \item \textbf{Score:} error of $\log(\text{rain} + 0.1\,\text{mm/h})$,
        always checked against mm bias, wet fraction and maximum. 3 seeds.
\end{itemize}
\end{frame}

%----------------------------------------------------------------------
\begin{frame}{How we cut the problem: a lessons-learned paper}
\begin{itemize}
  \item No method gives a real gain. A method paper is neither possible
        nor the point.
  \item What we have is general: how to tell in advance whether UDA can
        help, and how a failing method can look successful.
  \item Venue: AIES, \emph{Lessons Learned} article type. A full article
        only if COSMO and RainShift contrast cleanly (decided with Tom in
        December).
\end{itemize}
\medskip
\textbf{The five lessons}
\begin{enumerate}
  \item Measure whether $P(Y\mid X)$ shifts before adapting.
  \item Bound the problem with labels first (oracle, few-shot).
  \item Check that the adaptation term acts as intended.
  \item Evaluation choices can manufacture recovery.
  \item A generative model does not change the diagnosis.
\end{enumerate}
\end{frame}

%----------------------------------------------------------------------
\begin{frame}{Lesson 1: the relation differs between RainShift regions}
\small
Test: within bins of two inputs, compare the rain distributions of two
regions, against a null that shuffles whole weeks.

\medskip
\begin{tabular}{@{}lcc@{}}
\toprule
Pair & binned on (CAPE, water vapour) & binned on (ERA5 rain, CAPE) \\
\midrule
Europe -- Horn      & \yes{} (0.28 vs.\ null 0.03) & \yes{} (0.10 vs.\ 0.04) \\
Europe -- Melanesia & \yes{}, limited overlap      & \yes{} (0.06 vs.\ 0.02) \\
Horn -- Melanesia   & \no{} (0.07 vs.\ 0.07)       & \yes{} (0.10 vs.\ 0.02) \\
\bottomrule
\end{tabular}

\medskip
\begin{itemize}
  \item In COSMO, given the coarse rain, the fine-scale relation is nearly
        the same in every pair (difference $\leq 0.01$).
  \item Likely reason for the contrast: on RainShift, ERA5 rain and IMERG
        are different products, and how they relate differs by region.
  \item Caveat: upper bounds from two binned inputs; one calibration check
        is still to run.
\end{itemize}
\end{frame}

%----------------------------------------------------------------------
\begin{frame}{Lesson 2: labels first}
\begin{itemize}
  \item \textbf{Oracle} (trained with labels from both regions): into the
        tropics it reaches the target's own model, so the gap is fully
        recoverable in principle. Into Europe it leaves 22--39\,\% of the
        gap.
  \item \textbf{Few-shot}: about 2{,}000 labelled target hours bring a
        source model to the oracle in every direction.
  \item So the gap is reducible, and cheaply with labels, but unlabelled
        inputs do not tell the model which way to go.
\end{itemize}
\take{Before a UDA study, train the oracle and a few-shot curve. If a few
thousand labels close the gap, the practical question is getting labels,
not adaptation.}
\end{frame}

%----------------------------------------------------------------------
\begin{frame}{Lesson 3: the alignment did not act; BatchNorm did}
\begin{itemize}
  \item Every method that runs target data through the network also mixes
        target statistics into the BatchNorm layers.
  \item Re-estimate those statistics on source data only, weights
        unchanged: every method scores like no adaptation (U-Net, 3 seeds;
        AFM). One exception, and it gets there with absurd pixels.
  \item In the deployed network no method made the regions look alike: a
        classifier still tells them apart 98--100\,\% of the time.
  \item When we made the alignment act (BatchNorm kept out of the target
        pass), MMD reached the oracle's log score into Melanesia in 3 of 3
        seeds, but its mm error was worse in 2 of 3 (one seed up to
        $9\times10^{6}$\,mm/h).
  \item \textbf{Next check:} does plain ``more rain everywhere'' give the
        same log score?
\end{itemize}
\end{frame}

%----------------------------------------------------------------------
\begin{frame}{Lesson 4: evaluation can manufacture recovery}
\begin{itemize}
  \item Tropical models applied to Europe predict no rain at all. They
        score exactly like an all-zero field, so ``recovery'' there means
        leaving zero.
  \item Log scores hide absurd pixels: a DANN model with
        $10^{6}$\,mm/h pixels scored well.
  \item Horn and West Africa: the in-domain model is only 5--7\,\% better
        than predicting zero. No method has room to show anything there.
  \item The method ranking looked stable across seeds, but that came from
        BatchNorm. Without it, no ranking is stable.
  \item Every model, in-domain included, is far too dry.
\end{itemize}
\take{Always report: an all-zero and a climatology reference, mm bias, wet
fraction, maximum, and at least three seeds.}
\end{frame}

%----------------------------------------------------------------------
\begin{frame}{Lesson 5: a generative model does not change the diagnosis}
\begin{itemize}
  \item AFM's transfer gaps have the same shape as the U-Net's.
  \item Tropical sources into Europe: the ensemble collapses to nearly dry
        members.
  \item One partial gain in the whole project: FDA into Europe beats
        climatology in mm (CRPS 0.085--0.087 vs.\ 0.093\,mm). One model, two seeds.
  \item A probabilistic score can reward spread alone when the baseline
        ensemble has collapsed.
\end{itemize}
\take{Open: AFM still starts from a U-Net mean. Plain flow matching (CFM)
on COSMO tests whether the conclusion depends on that.}
\end{frame}

%----------------------------------------------------------------------
\begin{frame}{COSMO perfect model: first results (no adaptation yet)}
\begin{itemize}
  \item \textbf{Transfer mostly works:} 10 of 20 directions are within seed
        noise of the target's own model. On RainShift, none is.
  \item \textbf{Failures depend on direction:} dry source $\to$ wet target
        costs $+200$ to $+330\,\%$; the reverse costs $-3$ to $+6\,\%$.
  \item \textbf{The cost is missing rain occurrence:} failing models rain
        too rarely (wet fraction 0.05--0.09 vs.\ 0.11--0.14).
  \item \textbf{Seeds disagree in the failing cells:} the source data do
        not constrain the model there.
\end{itemize}
\take{Working hypothesis: a wet region contains the dry region's
conditions, not the reverse. Unlabelled target inputs can say
\emph{where} a model will fail, but not \emph{what} the rain is there.}
\end{frame}

%----------------------------------------------------------------------
\begin{frame}{Where each question stands}
\small
\renewcommand{\arraystretch}{1.25}
\begin{tabular}{@{}>{\raggedright\arraybackslash}p{0.34\linewidth}>{\raggedright\arraybackslash}p{0.30\linewidth}>{\raggedright\arraybackslash}p{0.30\linewidth}@{}}
\toprule
Question & RainShift & COSMO \\
\midrule
1.\ Does $P(Y\mid X)$ shift? & yes; Horn--Melanesia depends on the test & barely, given coarse rain \\
2.\ Recoverable, identifiable? & yes / no & oracle not run yet \\
3.\ Do methods act? & no; BatchNorm did & not run yet \\
4.\ Depends on the model? & no (AFM) & U-Net / AFM / CFM planned \\
5.\ Perfect model different? & --- & transfer mostly works \\
\bottomrule
\end{tabular}
\end{frame}

%----------------------------------------------------------------------
\begin{frame}{Next steps}
\textbf{Until the decision meeting (mid-December)}
\begin{itemize}
  \item Run the ``more rain everywhere'' check; finish lesson 3.
  \item Calibrate the lesson-1 test; finish the five-region RainShift
        matrix and its shift tests.
  \item A two-variable toy example of why unlabelled data cannot pick the
        right model.
  \item Spatial-structure scores (FSS, power spectra) for both datasets.
  \item Prepare the generative models on COSMO: test run, learning rate,
        ensemble scoring; then train AFM and CFM on 5 regions $\times$ 3
        seeds.
\end{itemize}
\textbf{January--February: COSMO experiments}
\begin{itemize}
  \item U-Net, AFM, CFM on 6 core directions (Spain$\leftrightarrow$Scandinavia,
        Poland$\leftrightarrow$British Isles, Alps$\leftrightarrow$British Isles):
        oracle, FDA, AdaBN, MMD; 3 seeds.
  \item U-Net on all 20 directions; few-shot; climate-change shift (+3\,K).
\end{itemize}
\end{frame}

%----------------------------------------------------------------------
\begin{frame}{Schedule}
\small


\begin{tabular}{@{}l|c|c|c|c|c@{}}
 & Nov & Dec & Jan & Feb & Mar \\
\hline
RainShift closed, text of lessons 1--5      & \cdone & \cplan &    &    &    \\
COSMO generative models: preparation, training & & \cplan &    &    &    \\
\textbf{Decision meeting with Tom} (18 Dec) &    & \cplan &    &    &    \\
COSMO: core directions, all three models    &    &    & \cplan &    &    \\
COSMO: U-Net on all directions, +3\,K shift &    &    & \cplan & \cplan &    \\
Final tables and figures, numbers frozen    &    &    &    & \cplan &    \\
Writing                                     &    &    &    &    & \cplan \\
\end{tabular}

\medskip
\begin{itemize}
  \item Experiments done by end of February 2027, writing in March.
  \item Compute: about 1{,}250 A100-hours for the core plan (about 9 days
        of our full GPU quota), about 2{,}800 for everything. Needs Tom's
        approval. Fallback: U-Net and CFM only.
\end{itemize}
\end{frame}

%----------------------------------------------------------------------
\begin{frame}{Decisions for the December meeting}
\begin{enumerate}
  \item \textbf{Paper type:} Lessons Learned (default) or a full article.
  \item \textbf{Compute} for the three-model comparison on COSMO: core
        only, full, or the reduced fallback.
  \item \textbf{How to state Horn--Melanesia} in lesson 1: depends on one
        test and on the calibration.
  \item \textbf{Title.}
\end{enumerate}
\end{frame}

\end{document}
